\documentclass[11pt,a4paper]{article}

% --- encoding & fonts ---
\usepackage[utf8]{inputenc}
\usepackage[T1]{fontenc}
\usepackage{lmodern}

% --- layout ---
\usepackage[margin=1in]{geometry}
\usepackage{microtype}
\setlength{\parskip}{0.4em}
\setlength{\parindent}{0pt}

% --- math & symbols ---
\usepackage{amsmath,amssymb}

% --- tables ---
\usepackage{booktabs}
\usepackage{array}

% --- graphics & code ---
\usepackage{graphicx}
\graphicspath{{figures/}}
\usepackage{xcolor}
\usepackage{listings}
\lstset{basicstyle=\ttfamily\footnotesize,breaklines=true,frame=single,columns=fullflexible,
        keepspaces=true,showstringspaces=false}

% --- bibliography (natbib + BibTeX; portable on Overleaf) ---
\usepackage[numbers,sort&compress]{natbib}
\bibliographystyle{unsrtnat}

% --- links (load hyperref last) ---
\usepackage{hyperref}
\hypersetup{colorlinks=true, linkcolor=black, citecolor=blue, urlcolor=blue,
            pdftitle={Gap2Idea}, pdfauthor={Yazan Alnakri; Hamza}}

% TODO: fill in full author names, affiliation, and supervisor.
\title{\textbf{Gap2Idea}: A Cost-Efficient End-to-End System for Mining\\
Research Gaps and Generating Evaluated Research Ideas}
\author{Yazan Alnakri \and Hamza}
\date{\today}

\begin{document}
\maketitle

\begin{abstract}
The volume of scientific literature makes it infeasible for researchers to
manually identify open research gaps and synthesise new directions. Large
language models (LLMs) can help, but applying them per paper does not scale:
extracting gaps from a million papers with a per-paper LLM costs on the order of
thousands of dollars and is rate-limited. We present \textbf{Gap2Idea}, an
end-to-end, deployable system that mines research gaps from scientific papers and
generates evaluated research ideas. Extraction uses a \emph{cost-efficient
funnel}: a free structural slice and a lightweight classifier reduce each paper
to a handful of candidate gap sentences, over which an LLM is applied only as a
precision filter, cutting cost by roughly two orders of magnitude ($\sim$\$3--31
vs $\sim$\$4{,}000 per million papers) while matching the quality range of a
fine-tuned transformer on a shared benchmark. We show empirically that extraction
quality is \emph{bounded by training data, not model capacity}, and adopt a
recipe that harvests mandated ``Limitations'' sections as clean supervision.
Extracted gaps are organised into a \emph{multi-relational gap-graph} whose
community bridges and frontier nodes seed \emph{novelty-by-recombination} idea
generation, performed by a \emph{grounded multi-agent} process with explicit
anti-hallucination constraints. Generated ideas are assessed by a multi-faceted
protocol combining a cross-provider LLM judge panel, automated novelty checks
against Semantic Scholar, and a human expert study. The system is open-source,
containerised, and deployed with a Model Context Protocol interface.
\emph{[Results sentence to be completed after the human study / graph-vs-random
ablation.]}
\end{abstract}

\tableofcontents
\newpage

\section{Research Questions}
\label{sec:rq}

Scientific output grows faster than any researcher can read. Turning that
literature into \emph{what has not yet been done}---concrete, novel, feasible
research directions---is a core scholarly task that does not scale manually, and
that na\"ive per-paper LLM use cannot scale \emph{economically}: a per-paper
extractor costs on the order of \$4{,}000 per million papers and is rate-limited.
This thesis builds and evaluates \textbf{Gap2Idea}, an end-to-end system that
mines research gaps from papers, organises them into a graph, and generates
evaluated research ideas. Three research questions structure both the system and
the literature; the pipeline instantiates them in sequence---\textbf{RQ2
(extract) $\rightarrow$ RQ1 (graph) $\rightarrow$ RQ3 (ideate)}.

\subsection{RQ2 --- Text-Based Gap Mining}
\emph{Can research gaps (limitations, future-work directions) be extracted from
scientific papers at corpus scale without a per-paper LLM, at a quality
competitive with fine-tuned transformers?}
\begin{itemize}
  \item \textbf{RQ2a.} What architecture localises and classifies gap sentences
        cheaply and at high recall? (A structural slice $\rightarrow$ a cheap
        classifier $\rightarrow$ an LLM precision filter.)
  \item \textbf{RQ2b.} Is extraction quality bounded by \emph{training data} or
        \emph{model capacity}?
  \item \textbf{RQ2c.} How does the funnel compare, on public benchmarks (LimGen,
        BAGELS, the future-work corpus), to fine-tuned and generative baselines?
\end{itemize}

\subsection{RQ1 --- Graph-Based Forecasting}
\emph{Can research gaps be organised into a multi-relational graph whose structure
identifies promising, novel directions---by recombination and bridging---better
than ungrounded pairing?}
\begin{itemize}
  \item \textbf{RQ1a.} How to embed and cluster gaps into interpretable
        communities?
  \item \textbf{RQ1b.} Do graph-derived \emph{bridge} and \emph{frontier} seeds
        yield better idea candidates than random gap pairs? (The extrinsic
        clustering signal.)
\end{itemize}

\subsection{RQ3 --- LLM-Driven Ideation}
\emph{Can an LLM, grounded in retrieved gap evidence, generate novel, feasible,
non-hallucinated research ideas, and can idea quality be evaluated reliably?}
\begin{itemize}
  \item \textbf{RQ3a.} What grounding and anti-hallucination constraints keep
        generated ideas faithful to evidence?
  \item \textbf{RQ3b.} Does a multi-agent critic/revise loop improve idea quality
        over a single pass?
  \item \textbf{RQ3c.} Can idea quality be assessed by combining an LLM-judge
        panel with a human expert study, and do LLM judgments correlate with human
        ratings?
\end{itemize}

\section{Literature Review}
\label{sec:lr}

We organise related work by the three research questions. Extraction-specific
claims are drawn from an adversarially-verified sweep (independent verifiers;
confidence noted in text where relevant); the graph and ideation clusters are
drawn from the reviewed corpus.

\subsection{RQ2 --- Text-based gap mining}
\label{sec:lr-rq2}

\paragraph{The task and its two shapes.} Work on ``gaps'' divides into
\emph{extraction} (locate the sentences stating a limitation or future-work
direction) and \emph{generation} (synthesise a limitations/future-work passage).
Extraction systems converge on a single architecture---\textbf{high-recall
structural/lexical retrieval $\rightarrow$ high-precision classifier
filter}---which is precisely the funnel this thesis uses. \citet{huwan2015}
pair a regex future-work extractor with a four-way classifier; \citet{zhangfws2022}
use two stages (a Naive-Bayes binary recogniser reaching Macro-F1 $90.7\%$ on
their corpus, then a SciBERT six-way typer at weighted-F1 $72.6\%$) and release a
labelled NLP/ACL future-work corpus (9{,}009 future-work vs 55{,}887 non-future-work
sentences). The closest analogue to our funnel is the randomized-controlled-trial
limitations system of \citet{rctlimits}: a keyword slice of candidate sentences
followed by a fine-tuned PubMedBERT filter reaching detection F1 $0.821$, scaled
to $\sim$12k articles with no per-paper LLM, and only marginally above a
rule-only baseline ($0.800$)---direct support for our design and a realistic
target. \citet{futuregen2025} state the pattern outright (regex high-recall
$\rightarrow$ LLM filter $\rightarrow$ RAG generation) and trim their input ``to
reduce API cost,'' i.e.\ they pay exactly the per-paper cost our cheap Stage~B is
designed to avoid.

\paragraph{Limitation datasets and generation.} \citet{limgen2024} harvest the
\emph{mandated} ``Limitations'' sections of 4{,}068 ACL papers as gold and
benchmark generative models---the mandated-section harvest is the
distant-supervision recipe we adopt (\S\ref{sec:method}). \citet{bagels} build a
limitations dataset over ACL/NeurIPS/PeerJ (verbatim author-stated plus
peer-review-supplemented) and benchmark generation by BERTScore/ROUGE and
coverage; its ACL split is the benchmark used in \S\ref{sec:results}. The same
group iterates on limitation \emph{generation}: topic-modelled aggregation
\citep{limtopic}, a per-paper graph with RAG \citep{generatingsuggesti2024}, and a
multimodal chart-limitation variant \citep{mitigatingvisual2024}. Sentence-level
classification of literature-review, problem, and method sentences is studied by
\citet{modellingclassifyi2025} and \citet{extractingproblem2026}; the latter
reports that fine-tuned models beat LLM in-context learning on this task and
diagnoses formulaic-expression shortcut-learning---a failure mode any cue-based
classifier inherits. \citet{gapmap2025} and \citet{canllms2025} push toward
\emph{implicit} gaps (not author-stated)---the harder, more useful frontier, and
the boundary of what a faithful \emph{extractor} (this thesis) deliberately does
not attempt. \citet{mixtureknowledge2024} bridge gap identification and graph
structure for review generation.

\paragraph{The hard problems, corroborated.} The literature independently
confirms our findings: gaps concentrate in \emph{terminal} sections
(Discussion/Future-Work/Conclusion) but some hide mid-paper; authors conflate
future-work with limitations, and $>$50\% of stated limitations cannot be
discerned from text alone; and while binary detection is easy, fine-grained
typing is hard for everyone (SciBERT wF1 $72.6\%$ on six future-work types;
PubMedBERT F1 $0.49$ on 24 limitation sub-types). Extraction datasets are
domain-locked (ACL, biomedical); none targets cross-domain arXiv AI/ML/math,
motivating our own small gold set.

\subsection{RQ1 --- Graph-based forecasting}
\label{sec:lr-rq1}

The reviewed RQ1 cluster frames ``finding promising directions'' as
graph/link-prediction-over-time, which motivates our gap-graph.

\paragraph{Knowledge-graph link prediction for trend forecasting.}
\citet{predictingfuture2022} (Science4Cast) forecast future concept links in a
100k-paper network; \citet{impact4cast2024} and \citet{forecastinghigh2024}
predict \emph{high-impact} future links on evolving knowledge graphs;
\citet{forecastingfuture2023} and the dual-link-prediction technology-opportunity
analysis of \citet{technologyopportun2023} do the same for AI and patents.
\citet{predictingnew2025} extract concept graphs with LLMs to predict new
materials-science directions.

\paragraph{Emerging-topic detection over dynamic graphs.}
\citet{trackingdynamics2020} track co-word-network dynamics for emerging topics;
\citet{studypredictabilit2025} predict new scholar topics from knowledge networks;
\citet{empiricalstudy} couple a knowledge graph with an LLM to reconstruct
scientific history and forecast trends.

\paragraph{Graph topic modelling and representation learning} supply the
machinery for our community step: graph-isomorphism and neural graph topic models
\citep{ginopic2024,ngtm2025,graphconvolutional2023}, graph-based topic-diversity
and multiview online clustering \citep{noveltopic2023,multiviewdeep2022}, and
self-supervised graph autoencoders and negative sampling
\citep{s2gae2023,understandingnegat2020}. \citet{linkprediction2024} and
\citet{literaturebased2025} tie graph structure directly to hypothesis
generation---the RQ1$\rightarrow$RQ3 bridge. Gap2Idea's contribution relative to
this cluster is to run the same graph/recombination machinery over
\emph{extracted research-gap nodes} (not concept or citation nodes) and to expose
\emph{bridge} (edge-betweenness) and \emph{frontier} structure as idea seeds.

\subsection{RQ3 --- LLM-driven ideation}
\label{sec:lr-rq3}

\paragraph{Knowledge-graph-grounded ideation.} Closest to our graph-seeded
generation with a human study is \citet{generationhuman2024}. \citet{chimera2025}
build a knowledge base of \emph{idea recombinations}; \citet{cam2023} mine
creative analogies; and \citet{ramonllull} realise combinatorial
theme$\times$domain$\times$method ideation---the explicit recombination
hypothesis.

\paragraph{Idea-development tools and grounding studies.} IdeaSynth
\citep{ideasynth2025}, SCI-IDEA \citep{sciidea2025}, workflow scaffolding
\citep{scaffoldingflexibl2025}, and CRISP-IM \citep{adaptingcrisp2021} evolve idea
facets with a human in the loop. Quality studies motivate our evidence-grounding
and human evaluation: \citet{canchatgpt2024} and \citet{comparingideation2024}
compare LLM and human ideation, and \citet{improvingresearch2025} report gains in
feasibility and quality from metadata grounding.

\paragraph{Surveys and role framings.} \citet{reviewllm2025} review 61
LLM-ideation studies; \citet{evolvingrole}, \citet{ai4research2025}, and
\citet{aiscience2025} frame LLMs as evaluator/collaborator/scientist and chart the
move toward agentic science. LLM-as-judge recurs across the cluster (including
ablation meta-evaluation, \citealp{abgen2025}); its standing caveat---%
self-evaluation bias when judge and generator share a model family---directly
motivates our cross-provider panel and human study.

\paragraph{Where Gap2Idea sits.} Individually, each component has precedent. The
contribution is their integration: a cost-efficient extraction funnel feeding a
gap-graph whose bridges and frontiers seed grounded, hallucination-gated,
multi-agent ideation, validated by a human-anchored evaluation
methodology---a deployed, reproducible pipeline rather than any single new model.

\section{Method}
\label{sec:method}

\subsection{System overview}
Six stages pass plain files (\texttt{data/} in, \texttt{artifacts/} out) and are
exposed as a CLI, a Streamlit app, and a Model Context Protocol (MCP) server:

\begin{lstlisting}
select-papers -> download-pdfs -> extract-text (GROBID; PyMuPDF fallback)
   -> gap extraction (funnel: Stage A/B/C)                     [RQ2]
   -> theme-mine (embed -> gap-graph: Leiden + bridge/frontier) [RQ1]
   -> generate-ideas (grounded, multi-agent, gated)            [RQ3]
   -> evaluate-ideas (LLM panel + S2 novelty + human study)
   -> export-ideas (LaTeX / PDF / full-paper)   |   serve-mcp
\end{lstlisting}

An LLM-provider abstraction (\texttt{pipeline/llm.py}) routes every LLM call
through one OpenAI-compatible client; \texttt{LLM\_PROVIDER} selects OpenRouter
(multi-model) or YandexGPT, with model-slug rewriting so that no call-site changes
when switching providers.

\subsection{Ingestion (RQ2 input)}
\texttt{extract-text} parses PDFs with GROBID \citep{grobid} (a TEI section tree;
ML-based, reliable on scrambled two-column PDFs), falling back to PyMuPDF when
GROBID is unavailable. GROBID ingestion is the corpus-scale ``PDF-parsing line
item'' and aligns with the AllenAI S2ORC pipeline, which also wraps GROBID
\citep{s2orc}; it feeds clean text and a real section tree to the whole pipeline
(\texttt{pipeline/grobid\_sections.py}, \texttt{gap\_funnel.slice\_grobid\_regions}).

\subsection{Gap extraction --- the funnel (RQ2)}
\label{sec:method-funnel}
Three stages of increasing cost, each shrinking the input for the next
(\texttt{pipeline/gap\_funnel.py}, \texttt{pipeline/gap\_llm\_filter.py}):

\begin{itemize}
  \item \textbf{Stage A --- structural slice} (\texttt{slice\_terminal\_regions} /
        \texttt{slice\_grobid\_regions}): on the reading-order sentence stream,
        union three recall sources---a span after a Limitations/Future-Work/%
        Conclusion \emph{heading}, a $\pm$window around an \emph{inline keyword}
        (catching column-scrambled headings), and the \emph{terminal tail}. With
        GROBID, keep Limitations/Future-Work/Discussion/Introduction sections and
        blacklist Related-Work/Background. Position-gated, priority-capped; free
        (regex/CPU); \textbf{drops $\sim$82\% of sentences}.
  \item \textbf{Stage B --- classify} (\texttt{cue\_label} +
        \texttt{EmbeddingGapHead}): high-precision cue rules give a free
        fast-accept and type; a logistic head over frozen \texttt{bge-small-en-v1.5}
        embeddings \citep{bgesmall} catches cue-less gaps, but only inside
        \emph{explicit} Limitations/Future-Work regions (tail/discussion use rules
        only, else false positives flood). Self-distilled from teacher labels with
        negatives drawn from the body \emph{outside} the slice; the shipped head
        adds $\sim$1{,}500 harvested ACL Limitations sentences as clean positives.
  \item \textbf{Stage C --- LLM precision filter}
        (\texttt{gap\_llm\_filter.LLMGapFilter}): an LLM judges only the $\sim$6
        survivors per paper and drops false positives (acknowledgments, formulas,
        citations, contribution claims, \emph{prior-work critiques}, PDF scramble).
        Two refinements introduced in this work: (i) a \emph{batched judge}---one
        structured \texttt{json\_schema} call per $\sim$40 candidates instead of
        one per sentence; and (ii) \emph{section-aware protection}---cue-rule hits
        inside explicit Limitations/Future-Work sections are trusted, while hits in
        Discussion/Intro/tail (where prior-work critiques concentrate) are judged.
\end{itemize}

A \emph{context ablation} found that feeding the judge the surrounding paragraph
($\pm 1$ sentence) or the paper title \emph{lowered} precision: a gap sentence and
its adjacent motivating/prior-work sentence have opposite ownership, so any local
window blurs the own-vs-prior boundary that the bare sentence draws cleanly.

\subsection{Gap graph --- theme mining (RQ1)}
\texttt{pipeline/theme\_mining.py} and \texttt{gap\_graph.py}: gap sentences are
embedded, a multi-relational graph is built (semantic-similarity, same-paper,
same-section, and shared-method edges), Leiden community detection
\citep{leiden} partitions it, and edge-betweenness identifies \emph{bridge} gaps
(spanning communities) and \emph{frontier} nodes; these seed
novelty-by-recombination. \emph{(Method verified against code; extrinsic
evaluation pending, \S\ref{sec:progress}.)}

\subsection{Idea generation (RQ3)}
\texttt{pipeline/openai\_ideas.py}, \texttt{agents.py}, \texttt{orchestrator.py}:
five modes (\textbf{bridge}, \textbf{within-community}, \textbf{method-gap},
\textbf{frontier}, \textbf{orchestrated}) turn seeds into ideas, grounded in
retrieved evidence with anti-hallucination gates---a verbatim-evidence
constraint, an evidence-overlap check, and a required named baseline plus
falsifiable prediction. The orchestrated mode runs a multi-agent
critic/revise/sanity loop (\texttt{agents.py}, \texttt{sanity.py}).

\subsection{Evaluation (cross-cutting)}
\texttt{pipeline/evaluation.py}: a cross-provider LLM-judge panel scores ideas
(novelty / specificity / feasibility / grounding) with inter-judge agreement; an
automated Semantic-Scholar novelty check (\texttt{semantic\_scholar.py}); and a
human expert study form (Krippendorff's $\alpha$), with LLM-judge-vs-human
correlation as the justification for using the panel at all.

\subsection{Output and deployment}
\texttt{export.py}, \texttt{paper\_drafter.py}: per-idea LaTeX
(\texttt{minimal}/\texttt{standard}/\texttt{ieee}), a consolidated PDF, or a
full-paper plan; \texttt{mcp\_server.py} exposes the corpus to Claude Desktop and
Cursor. The system is containerised (Docker), deployable on Cloud Run, and covered
by CI and a test suite.

\section{Results}
\label{sec:results}

Extraction (RQ2) is the benchmarked component. RQ1/RQ3 results are pending real
runs (\S\ref{sec:progress}) and are reported honestly, not fabricated.

\subsection{Stage A --- localisation (the recall ceiling)}
On the clean 19-gap / 9-paper gold set (token-containment match at threshold
$\tau$):

\begin{table}[h]
\centering
\begin{tabular}{@{}cccc@{}}
\toprule
containment $\tau$ & recall (all) & future-work & limitation \\
\midrule
$0.80$ & \textbf{0.84} & \textbf{0.91} & 0.75 \\
$0.70$ & 0.89 & 1.00 & 0.75 \\
\bottomrule
\end{tabular}
\end{table}

Stage A drops $\sim$82\% of sentences for free. Future-work localisation is
effectively solved; the residual misses are mid-paper own-work limitations
(a structural, not a slicer, limit).

\subsection{Stage B / end-to-end, and the data-not-model finding}

\begin{table}[h]
\centering
\begin{tabular}{@{}lcccc@{}}
\toprule
Stage B head & preds/paper & gap recall & limitation recall & type acc \\
\midrule
rules only & 2.2 & 0.32 & 0.11 & 1.00 \\
hybrid (bge+logreg) & 4.2 & 0.42 & 0.11 & 1.00 \\
\textbf{+ ACL limitations} \emph{(shipped)} & 6.1 & \textbf{0.53} & \textbf{0.44} & 0.90 \\
\bottomrule
\end{tabular}
\end{table}

\textbf{Data, not model} (a clean negative result): classifier method
(logreg $\approx$ DistilBERT $\approx$ SetFit) \emph{and} frozen encoder
(bge-small $\approx$ bge-base $\approx$ mpnet $\approx$ SPECTER \citep{specter})
all tie at limitation recall $\approx 0.11$; only \emph{clean data}---harvesting
mandated ACL Limitations sections \citep{limgen2024}---moved limitation recall
$0.11 \rightarrow 0.44$ ($4\times$) and end-to-end $0.42 \rightarrow 0.53$,
saturating at $\sim$1{,}500 sentences. Upgrading the model is a dead end, and the
literature agrees (even a domain-pretrained BERT caps $\sim$0.5 on limitation
typing).

\subsection{Cost (RQ2's headline)}

\begin{table}[h]
\centering
\begin{tabular}{@{}lc@{}}
\toprule
Approach & Cost / 1M papers \\
\midrule
Per-paper LLM (gpt-4.1-mini, $\sim$9k tokens) & $\sim$\$4{,}000 \\
Funnel (Stage A slice + Stage B head) & $\sim$\$3--31 \\
Stage C batched judge (this work) & $\sim$\$25--30 \\
\bottomrule
\end{tabular}
\end{table}

The funnel is $\sim$130--160$\times$ cheaper; adding the batched Stage C keeps it
within the same order ($\sim$\$30/1M vs $\sim$\$180/1M for a per-sentence LLM
filter---a $-88\%$ token reduction from batching).

\subsection{Same-data comparison vs prior art (LimGen)}
Leakage-clean, on LimGen ACL data, our cheap frozen \emph{stacking ensemble} F1
$0.627$ beats a reproduced Bernoulli-NB \citep{zhangfws2022} ($0.553$), TF-IDF
($0.610$), and base bge ($0.610$), and is $\sim$0.05 behind a fine-tuned
DistilBERT ($0.674$). On a 196-sentence sample the Stage~C lift is
$0.643 \rightarrow 0.725$ (precision $0.56 \rightarrow 0.79$)---\emph{indicative}
(given the sample size) that the funnel matches the range of a fine-tuned
transformer at a fraction of the cost.

\subsection{New benchmarks (this work)}

\paragraph{BAGELS extraction (ACL).} Over 127 papers / 1{,}212 verbatim gold
limitation sentences \citep{bagels}, the funnel's Stage~A+B coverage of the gold
is \textbf{0.738 verbatim} and \textbf{0.790 semantic} (bge cosine $\geq 0.75$);
per table, ACL-23 is $0.822/0.859$ and ACL-24 $0.699/0.758$. BAGELS is a
\emph{generation} benchmark whose best generation Coverage-of-Ground-Truth is
$76.62\%$ (GPT-3.5). On the comparable coverage metric our \emph{extraction}
matches or exceeds their best \emph{generation}---but this is an easier task
(the limitations section is present) with a stricter matcher, so it is framed as
\emph{establishing the first cheap extraction baseline on BAGELS}, not beating
their leaderboard. Stage~C is a precision filter; on a limitations-only gold set
(no negatives) it can only cost coverage, so A+B is the coverage figure to quote.

\paragraph{Future-work sentence corpus.} On the labelled corpus of
\citet{zhangfws2022} (64{,}896 sentences), sentence classification is a
lexical-cue task on pre-segmented sentences, which favours bag-of-words. Our head
scores \textbf{macro-F1 0.765 zero-shot} and \textbf{0.772 with Stage C} (Stage~C
removed 314 false positives against 118 true future-work sentences, raising
precision $0.80 \rightarrow 0.88$). We do \emph{not} beat their reported $0.907$;
however, that number is \emph{not reproducible}: a faithful reproduction of their
own Bernoulli-NB on the full corpus reaches $0.820$, and all text-only methods
(their NB, TF-IDF, our embeddings) cluster at $0.76$--$0.82$. These pre-segmented
sentence benchmarks do not exercise the funnel's actual contributions (structural
localisation and corpus-scale extraction from raw PDFs).

\paragraph{Provider robustness.} YandexGPT occasionally issues a content-filter
refusal on a batch; the per-sentence fallback in Stage~C absorbs it with no lost
paper---a small reliability data point for single-provider deployments.

\section{Current Progress}
\label{sec:progress}

Component scorecard (\emph{built} $\neq$ \emph{evaluated}):

\begin{table}[h]
\centering
\small
\begin{tabular}{@{}p{3.4cm}cp{4.2cm}p{3.6cm}@{}}
\toprule
\textbf{Component} & \textbf{Built} & \textbf{Evaluated} & \textbf{State} \\
\midrule
RQ2 --- Gap extraction (funnel A/B/C) & yes & small gold + LimGen + BAGELS/FWS (this work) & \textbf{paper-ready}, honest/indicative claims \\
RQ1 --- Gap graph (Leiden + bridge/frontier) & yes (novel) & prior bench was $N{=}11$ and excluded the shipped graph method; no extrinsic eval & \textbf{needs real eval} on hundreds of gaps \\
RQ3 --- Idea generation (5 modes + multi-agent) & yes & committed runs are the simple path; no orchestrated/critic runs saved & \textbf{needs runs} \\
Idea evaluation / novelty & yes (code) & panel unused (single judge, self-bias); human study has no real responses; S2 novelty on $\sim$7/27 & \textbf{make-or-break; not yet credible} \\
Output + deployment (MCP, Docker, CI) & yes (strong) & 19 tests & \textbf{paper-ready} \\
\bottomrule
\end{tabular}
\end{table}

\subsection{Done in this cycle (extraction hardening + benchmarking)}
\begin{itemize}
  \item YandexGPT provider integration behind the LLM-provider abstraction
        (model-slug rewriting; strict \texttt{json\_schema} supported).
  \item A \textbf{batched Stage C} judge ($\sim$2 calls per 25 papers,
        $-88\%$ tokens vs one call per sentence) with prior-work and
        self-promotion rejection and section-aware rule protection.
  \item A \textbf{context ablation}: feeding the judge the surrounding paragraph
        ($\pm 1$ sentence) or the paper title \emph{lowered} precision (bare
        sentence $>$ sentence-in-context), a reproducible negative result.
  \item The \textbf{BAGELS extraction benchmark} (0.74 verbatim / 0.79 semantic
        coverage), with a fair comparison to BAGELS's own generation numbers.
  \item The \textbf{future-work benchmark} plus a \textbf{reproducibility
        finding}: the reported SOTA of 0.907 does not reproduce; $\sim$0.82 is the
        reproducible bar, and all text-only methods cluster at 0.76--0.82.
  \item A \textbf{dataset-suitability survey} concluding that no ideal
        cross-domain, full-paper, verbatim extraction benchmark exists (BAGELS is
        the best available; Dr.\ Inventor's rhetorical layer is unavailable).
\end{itemize}

\subsection{Highest-value next steps (priority order)}
\begin{enumerate}
  \item \textbf{P0 --- evaluation integrity.} Remove any placeholder human data;
        run a \emph{real} human study (5--10 domain readers, $\sim$15 ideas,
        Krippendorff's $\alpha$); run the cross-provider judge panel and report
        \emph{LLM-vs-human correlation}; fix Semantic-Scholar novelty coverage.
  \item \textbf{P1 --- justify the pipeline logic.} The \emph{graph-vs-random
        ablation}: do bridge/frontier seeds produce ideas rated higher (panel and
        humans, blind)? This is the experiment that earns every upstream component
        its place. Plus a with/without-critic ablation for the multi-agent path.
  \item \textbf{P2 --- component evals at real size.} Re-run clustering on hundreds
        of gaps from \texttt{runs/*} \emph{including} the shipped Leiden-graph
        method; enlarge the extraction gold (19 $\rightarrow$ a few hundred).
\end{enumerate}

\paragraph{Framing.} The intended contribution is a \emph{system} and an
\emph{evaluation methodology}---cost-efficient extraction, graph-seeded
recombination, anti-hallucination engineering, deployability, and a
human-validated idea-quality study---rather than a ``beats SOTA'' claim.


\bibliography{references}

\end{document}
